\uuid{aW1E}
\chapitre{Statistique}
\niveau{L2}
\module{Analyse de données}
\sousChapitre{Tests d'hypothèses, intervalle de confiance}
\titre{Contrôle qualité}
\theme{statistiques, tests d'hypothèses}
\auteur{}
\datecreate{2022-10-17}
\organisation{AMSCC}
\difficulte{}
\contenu{
\texte{$\mathcal N(\mu,\sigma)$ désigne la loi normale de moyenne $\mu$ et d’écart-type $\sigma$. Les quantiles sont à gauche : $t_{\nu;\gamma}$ pour la loi $St(\nu)$, avec une probabilité $\gamma$ à gauche du quantile.}
\texte{Un pharmacien souhaite être livré avant 18 h. Les durées de 13 trajets du coursier, supposées indépendantes de loi $\mathcal N(\mu,\sigma)$, sont en minutes :
\[20,30,45,39,16,25,23,19,40,21,30,19,40.\]
Le pharmacien propose une moyenne de 25 minutes et le grossiste de 30 minutes.}
\begin{enumerate}
\item \question{Estimer la moyenne et l'écart-type ; rappeler les propriétés des estimateurs.}
\reponse{$\overline X$ et $S^2$ sont sans biais pour $\mu$ et $\sigma^2$. On obtient $\bar x=28{,}23077$, $s^2=96{,}52564$ et $s\simeq9{,}82475$. $S$ est la racine de la variance corrigée et n'est généralement pas sans biais pour $\sigma$.}
\item \question{Tester $H_0:\mu=25$ contre l'alternative $\mu=30$ au seuil de $5\%$.
\begin{enumerate}
\item Donner la statistique et sa loi sous $H_0$.
\item Déterminer la région critique.
\item Peut-on calculer un risque de seconde espèce numérique sans connaître $\sigma$ ? Exprimer sa dépendance en $\sigma$ ; pour illustration, prendre $\sigma=10$ minutes uniquement pour ce calcul de puissance.
\item Prendre une décision avec les données et expliquer sa portée.
\end{enumerate}}
\indication{L'alternative est supérieure à 25 : le test est unilatéral à droite. Pour la puissance exacte de Student, il faut la loi de Student non centrale.}
\reponse{La statistique $T=\frac{\overline X-25}{\frac{S}{\sqrt{13}}}$ suit $St(12)$ sous $H_0$. On rejette pour $t_{\mathrm{obs}}>t_{12;0{,}95}\simeq1{,}78229$.
Sous $\mu=30$, $T$ suit une Student non centrale à 12 ddl de paramètre $\delta=\frac{5\sqrt{13}}{\sigma}$. Le risque est
\[\beta(\sigma)=\mathbb P\bigl(T_{12,\delta}\leq1{,}78229\bigr).\]
Il n'a pas de valeur unique sans $\sigma$. À $\sigma=10$, un logiciel donne $\beta\simeq0{,}477988$ (puissance $0{,}522012$). Cette valeur de $\sigma$ sert à décrire une alternative, et ne remplace pas $S$ dans le test.
Avec les données, $t_{\mathrm{obs}}\simeq1{,}18565$ et $p_{\mathrm{val}}\simeq0{,}12935$ : non-rejet à $5\%$. On ne prouve pas une moyenne de 25 minutes. Une analyse de moyenne ne garantit pas non plus l'heure d'arrivée de chaque livraison.}
\end{enumerate}
}
